\documentclass[11pt]{article}
\usepackage[margin=1in]{geometry}
\usepackage{amsmath,amssymb,amsthm,hyperref}
\newtheorem{theorem}{Theorem}
\newtheorem{lemma}[theorem]{Lemma}
\newcommand{\E}{\mathbb E}
\newcommand{\R}{\mathbb R}
\title{A counterexample family to Simon's Le Roy-type criterion}
\author{Principia Math In Progress}
\date{11 September 2026}
\begin{document}
\maketitle
\begin{abstract}
We disprove the proposed necessity of the logarithmic infinite-divisibility
criterion for complete monotonicity of a Le Roy-type Mittag--Leffler
function. With $\alpha=1/3$ and $\gamma=2$, a nonempty open interval of
parameters $\beta$ below the proposed threshold still gives completely
monotone functions. The proof constructs a signed density by beta
integrals, proves strict positivity at the threshold, and shows that
positivity persists across that threshold. The existence of the interval
is proved analytically; its width is not assigned a numerical value.
\end{abstract}

\section*{Task 1: Resolve MathDB problem 349140}
Write
\[
 F_{\alpha,\beta}^{(\gamma)}(z)
 =\sum_{n=0}^\infty\frac{z^n}{\Gamma(\beta+\alpha n)^\gamma},
 \qquad \alpha,\beta,\gamma>0.
\]
The conjecture in Section 3 of \cite{source} asserts the equivalence
\begin{equation}\label{eq:conjecture}
 F_{\alpha,\beta}^{(\gamma)}(-x)\text{ is completely monotone for }x>0
 \quad\Longleftrightarrow\quad
 z+\gamma(z^{\beta-\alpha}-z^\beta)\leq1\quad(0<z<1).
\end{equation}
Its sufficient direction was already proved there. We refute necessity.

\begin{theorem}\label{thm:counter}
Define
\begin{equation}\label{eq:threshold}
 \beta_c=\max_{0<y<1}
 \frac13\left(1+\frac{\log((1+y+y^2)/2)}{\log y}\right).
\end{equation}
Then $1/3<\beta_c<1/2$, and there exists $\varepsilon>0$ such that,
for every $\beta\in(\beta_c-\varepsilon,\beta_c)$,
$F_{1/3,\beta}^{(2)}(-x)$ is strictly completely monotone on $(0,\infty)$,
whereas the right side of \eqref{eq:conjecture} fails.
\end{theorem}
The proof below defines a permissible positive $\varepsilon$ in terms of
convergent integrals and extrema on compact sets. It does not claim that
any particular decimal or rational $\beta$ has been certified.

\subsection*{The threshold}
Set $a=1/3$ and $b=2/3$ throughout. With $z=y^3$, the proposed condition is
equivalent to
\begin{equation}\label{eq:polynomial}
 1+y+y^2\geq2y^{3\beta-1}\quad(0<y<1),
\end{equation}
because the difference between its two sides in \eqref{eq:conjecture}
is $(1-y)(1+y+y^2-2y^{3\beta-1})$.
Taking logarithms shows that \eqref{eq:polynomial} holds exactly when
$\beta\geq\beta_c$.
The function being maximized tends to $a$ at zero and to $-\infty$ at one,
and is greater than $a$ whenever $1+y+y^2<2$. Its maximum is therefore
attained at some $y_*\in(0,1)$ and is greater than $a$.
Moreover,
\[
 1+y+y^2-2\sqrt y=(1-\sqrt y)^2+y^2>0,
\]
so every value in \eqref{eq:threshold}, including its attained maximum,
is less than $1/2$.
At any maximizer $y_*$, equality holds in \eqref{eq:polynomial} for
$\beta=\beta_c$. For every $\beta<\beta_c$, the same $y_*$ gives a strict
violation, with $z=y_*^3$ in \eqref{eq:conjecture}.

\subsection*{A continuous family of signed densities}
For $a<\beta<1/2$, define
\begin{align}
 H_\beta(s)&=\frac{\Gamma(s+a)\Gamma(s+b)\Gamma(s+1)}
                       {\Gamma(s+\beta)^2},\label{eq:mellin}\\
 w_{c,\beta}(u)&=\frac{u^{c-1}(1-u)^{\beta-c}}
                         {\Gamma(\beta+1-c)},
 \qquad c\in\{a,b\},\quad 0<u<1,\label{eq:weights}\\
 P_\beta(t)&=t^2+(2\beta-3)t+(\beta-1)^2,\label{eq:quadratic}\\
 h_\beta(x)&=\int_0^1\int_0^1
 w_{a,\beta}(u)w_{b,\beta}(v)
 \frac{e^{-x/(uv)}}{uv}\,P_\beta\!\left(\frac{x}{uv}\right)\,du\,dv,
 \qquad x>0.\label{eq:density}
\end{align}
Both beta weights are integrable and positive. The function $h_\beta$
may have either sign at this stage.

\begin{lemma}\label{lem:properties}
The family $h_\beta(x)$ is jointly continuous, with a continuous
$\beta$ derivative, on $(a,1/2)\times(0,\infty)$.
For every complex $s$ with $\Re s>-a$, it is absolutely integrable
against $x^s$ and satisfies
\begin{equation}\label{eq:transform}
 \int_0^\infty x^s h_\beta(x)\,dx=H_\beta(s).
\end{equation}
On each compact parameter interval $I\subset(a,1/2)$, it is positive
for all sufficiently small $x>0$, uniformly in $\beta\in I$.
It is also positive for every $x\geq3$ and every $\beta\in(a,1/2)$.
\end{lemma}
\begin{proof}
For $s$ in the stated half-plane, absolute Fubini is justified by
putting $t=x/(uv)$ in \eqref{eq:density}. The absolute $t$ integral is
bounded by a sum of gamma integrals, since $\Re s>-a>-1$.
The remaining beta integrals have exponents $s+a-1$ and $s+b-1$
at zero and integrable exponents at one. Furthermore,
\[
 \int_0^\infty t^s e^{-t}P_\beta(t)\,dt
 =(s+\beta)^2\Gamma(s+1),\qquad
 \int_0^1 u^s w_{c,\beta}(u)\,du
 =\frac{\Gamma(s+c)}{\Gamma(s+\beta+1)}.
\]
Their product is \eqref{eq:transform}. In particular $h_\beta\in L^1$.
On compact positive $x$ intervals, the exponential in
\eqref{eq:density} dominates all powers of $(uv)^{-1}$; logarithmic
factors from a $\beta$ derivative are integrable at $u=1,v=1$.
This proves the asserted continuity by dominated differentiation.

Here are details of the uniform small-$x$ assertion. Mellin inversion gives
\[
 h_\beta(x)=\frac1{2\pi i}\int_{-i\infty}^{i\infty}
                   H_\beta(s)x^{-s-1}\,ds.
\]
Indeed, \eqref{eq:transform} at $s=it$ is the Fourier transform of the
integrable continuous function $e^v h_\beta(e^v)$, and is itself integrable.
Uniform Stirling estimates on compact $\beta$ intervals give
\[
 |H_\beta(\sigma+it)|\leq C_I(1+|t|)^{A_I}e^{-\pi|t|/2}
 \quad(-1/2\leq\sigma\leq0,\ |t|\geq1).
\]
We may therefore move the contour to $\Re s=-1/2$. Only the simple pole
$s=-a$ is crossed. Consequently,
\begin{equation}\label{eq:small}
 h_\beta(x)=C_\beta x^{-2/3}+R_\beta(x),\qquad
 C_\beta=\frac{\Gamma(1/3)\Gamma(2/3)}{\Gamma(\beta-1/3)^2}>0,
 \qquad |R_\beta(x)|\leq M_I x^{-1/2},
\end{equation}
where one can take
\[
 M_I=\frac1{2\pi}\sup_{\beta\in I}
          \int_{\R}|H_\beta(-1/2+it)|\,dt<\infty.
\]
The minimum of $C_\beta$ on $I$ is positive, proving uniform positivity
near zero. Finally, for $t\geq3$ and $\beta>a$,
\[
 P_\beta(t)>t(t-7/3)>0.
\]
When $x\geq3$, the argument $t=x/(uv)$ in \eqref{eq:density} is at
least three. The entire integrand is then positive.
\end{proof}

\subsection*{Strict positivity at the threshold}
Nonnegativity alone at the threshold would not suffice for perturbation
on a compact $x$ interval. The following argument establishes the needed
pointwise strictness.

\begin{lemma}\label{lem:strict}
For every $x>0$, $h_{\beta_c}(x)>0$.
\end{lemma}
\begin{proof}
The standard exponential formula for the gamma function yields, for $s\geq0$,
\begin{equation}\label{eq:levy}
 \log\frac{H_\beta(s)}{H_\beta(0)}
 =d_\beta s+\int_0^\infty(e^{-su}-1+su)\frac{N_\beta(u)}u\,du,
\end{equation}
where
\[
 d_\beta=\psi(a)+\psi(b)+\psi(1)-2\psi(\beta),\qquad
 N_\beta(u)=\frac{e^{-au}+e^{-bu}+e^{-u}-2e^{-\beta u}}{1-e^{-u}}.
\]
This is the same gamma/L\'evy--Khintchine calculation used in
\cite[Section 2]{source}, with a change of scale.
At $\beta=\beta_c$, \eqref{eq:polynomial} gives $N_\beta\geq0$.
Moreover $N_\beta(u)/u\sim u^{-2}$ at zero and decays exponentially at
infinity. Thus \eqref{eq:levy} is the cumulant of an infinitely divisible
real random variable $T$, with negative jumps and all positive
exponential moments. The formula also holds on the imaginary axis.
By \eqref{eq:transform} and uniqueness of Fourier transforms of finite
signed measures, its continuous density is
\begin{equation}\label{eq:logdensity}
 g(t)=\frac{e^t h_{\beta_c}(e^t)}{H_{\beta_c}(0)}.
\end{equation}
In particular, this expression is nonnegative everywhere.

Since $a<\beta_c<b<1$, the term $e^{-au}$ dominates the numerator of
$N_{\beta_c}$ at infinity. Choose $L>0$ so that $N_{\beta_c}(u)>0$
for all $u\geq L$. Choose a continuous bounded function $\nu_1$ which
vanishes for $u\leq L$, is strictly positive for $u>L$, and satisfies
\[
 0\leq\nu_1(u)\leq N_{\beta_c}(u)/u.
\]
For example multiply that density on $u>L$ by
$1-e^{-(u-L)}$ and set it to zero elsewhere. Its mass
$\Lambda=\int\nu_1(u)\,du$ is finite and positive, and $\nu_1$ is
uniformly continuous.
Split the L\'evy measure to write $T=W+Z$, independently, where $Z$ is
the negative compound-Poisson variable with intensity density $\nu_1$.
The remaining drift is absorbed in $W$.

For $s\geq0$,
\[
 -\Lambda\leq\log\E e^{sZ}
 =\int_0^\infty(e^{-su}-1)\nu_1(u)\,du\leq0.
\]
Stirling's formula gives
$\log(H_{\beta_c}(s)/H_{\beta_c}(0))=s\log s+O(s)$.
It follows that $\log\E e^{sW}=s\log s+O(s)$, so $W$ has unbounded
support above: an upper bound $W\leq A$ would imply
$\log\E e^{sW}\leq As$.

The event that $Z$ has exactly one jump gives the continuous density
minorant
\[
 g_1(t)=e^{-\Lambda}\int_{\R}\nu_1(w-t)\,\mathbb P(W\in dw).
\]
For every real $t$, the unbounded upper support of $W$ and the positivity
of $\nu_1$ on $(L,\infty)$ imply $g_1(t)>0$.
Continuity follows from uniform continuity and boundedness of $\nu_1$.
The probability-measure domination gives $g\geq g_1$ almost everywhere;
as both are continuous, it holds everywhere.
Equation \eqref{eq:logdensity} now proves the lemma.
\end{proof}

\subsection*{Perturbation and the Mittag--Leffler transform}
Choose a closed interval $I\subset(a,1/2)$ with $\beta_c$ in its interior.
Let $c_I=\min_{\beta\in I} C_\beta>0$ and choose $\eta\in(0,1]$ so small
that $M_I\eta^{1/6}<c_I/2$. By \eqref{eq:small},
$h_\beta(x)>0$ for $0<x\leq\eta$, uniformly for $\beta\in I$.
Lemma \ref{lem:strict} gives
\[
 m=\min_{\eta\leq x\leq3}h_{\beta_c}(x)>0.
\]
Joint continuous differentiability from Lemma \ref{lem:properties} gives
\[
 A=\sup_{\beta\in I,\ \eta\leq x\leq3}
                |\partial_\beta h_\beta(x)|<\infty.
\]
Take any positive $\varepsilon$ smaller than both the distance from
$\beta_c$ to the endpoints of $I$ and $m/[2(1+A)]$.
For $|\beta-\beta_c|<\varepsilon$, the mean value theorem gives
$h_\beta(x)>m/2$ on $[\eta,3]$. Together with the uniform small-$x$
bound and positivity for $x\geq3$, this proves
\begin{equation}\label{eq:perturbed}
 h_\beta(x)>0\quad\text{for all }x>0,
 \qquad |\beta-\beta_c|<\varepsilon.
\end{equation}

Normalize $h_\beta/H_\beta(0)$ to be the density of a positive random
variable $Y_\beta$, and put $X_\beta=3Y_\beta^{1/3}$.
Gamma triplication and \eqref{eq:transform} give, for every $s\geq0$,
\begin{equation}\label{eq:moments}
 \E X_\beta^s
 =3^s\frac{H_\beta(s/3)}{H_\beta(0)}
 =\frac{\Gamma(1+s)\Gamma(\beta)^2}{\Gamma(\beta+s/3)^2}.
\end{equation}
The series for $F_{1/3,\beta}^{(2)}$ is entire, as follows directly from
Stirling's formula. For every real $z\geq0$, monotone convergence and
\eqref{eq:moments} therefore give
$\E e^{zX_\beta}=\Gamma(\beta)^2 F_{1/3,\beta}^{(2)}(z)<\infty$.
Absolute convergence then justifies the same identity for complex $z$.
In particular,
\begin{equation}\label{eq:laplace}
 F_{1/3,\beta}^{(2)}(-x)
 =\frac1{\Gamma(\beta)^2}\E e^{-xX_\beta}\quad(x>0).
\end{equation}
For every integer $k\geq0$, differentiation under this probability
integral is justified by boundedness of $t^k e^{-xt/2}$ locally in $x>0$.
Since $X_\beta>0$ almost surely, all alternating derivatives are strictly
positive. Finally choose $\beta\in(\beta_c-\varepsilon,\beta_c)$.
The threshold calculation shows failure of the proposed condition,
while \eqref{eq:laplace} proves complete monotonicity. This completes the
proof of Theorem \ref{thm:counter}.

\section*{Task 2: Certify an explicit counterexample parameter}
Replace the implicit positive width $\varepsilon$ by a rigorous numerical
lower bound. In particular, determine whether $\beta=43/100$ is a
counterexample, using certified bounds for the small-$x$ remainder and
the compact-interval minimum. A positive numerical grid alone is not
a certificate.

\section*{Task 3: Determine the actual complete-monotonicity threshold}
For $\alpha=1/3$ and $\gamma=2$, determine the infimum of the parameters
$\beta$ giving complete monotonicity. Describe the loss of positivity
of $h_\beta$ at this boundary, including whether it occurs through a
unique double zero at a finite positive argument.

\section*{Task 4: Extend the counterexamples to other parameters}
Establish an analogous stability theorem for other rational $\alpha$
and integer $\gamma$, and then for noninteger $\gamma$. Determine
whether the counterexamples occupy a nonempty open subset of the full
three-dimensional parameter space.

\section*{Task 5: Find a direct positive representation below the threshold}
Construct an explicit positive mixture or product representation of
$X_\beta$ for counterexample parameters. Such a representation should
make positivity visible without perturbing from $\beta_c$, and explain
how the sign-changing gamma L\'evy kernel is compatible with a positive
probability law.

\begin{thebibliography}{9}
\bibitem{source} T. Simon,
\emph{Remark on a Mittag--Leffler function of Le Roy type},
Integral Transforms Spec. Funct. \textbf{33} (2022), 108--114,
\href{https://doi.org/10.1080/10652469.2021.1913138}{doi:10.1080/10652469.2021.1913138}.
\href{https://arxiv.org/abs/2012.11018}{arXiv:2012.11018v2},
25 March 2021, Section 3, pages 5--6.
\end{thebibliography}
\end{document}
